% ex_cap5.tex
\chapter*{Considerações finais}
\addcontentsline{toc}{chapter}{Considerações finais}
\markboth{Considerações finais}{Considerações finais}

Em março de 2022, no início do trabalho de campo, João me guiou pelas instalações físicas do Centro. Ao me aproximar da sala principal de convivência dos pesquisadores, notei que a porta estava fechada, mas uma parede inteira de vidro expunha todo o interior ao corredor. Parei ali por alguns instantes, observando através do vidro. Fileiras de computadores ocupavam estações de trabalho compartilhadas, e naquele dia não havia ninguém na sala. A transparência seletiva daquele espaço me intrigou: a parede de vidro permitia ver o interior, mas a porta fechada controlava quem entrava. Será que era ali o \enquote{laboratório do \texttt{C4AI}}? Mais adiante, João me mostrou outra sala, esta com as máquinas, as \texttt{GPUs} refrigeradas e enfileiradas que processavam os dados e treinavam os modelos, e a apresentou como algo próximo da área de \texttt{TI} do Centro. Não havia cientistas entre elas, e o único movimento era o piscar rítmico de suas luzes. Perguntei se era ali o laboratório. João disse que não. Respondi com um \enquote{Ah!} monossilábico e seguimos caminhando.

Comecei esta tese procurando um laboratório e encontrei uma sala de vidro diante de uma sala de máquinas. A pergunta com que abri o capítulo~\ref{capitulo1}, onde está o laboratório dos pesquisadores de inteligência artificial e onde estão os seus cientistas, nasceu desse desencontro entre o que eu esperava ver e o que o campo me dava a ver. Eu procurava o equivalente, na pesquisa em IA, dos espaços que a antropologia da ciência aprendeu a descrever na física e na biologia, com suas bancadas, seus instrumentos e suas rotinas localizadas, e encontrei um \enquote{laboratório} que não correspondia a essa expectativa. Não encontrei os cientistas reunidos naquela sala, nem ali nem nos quatro anos seguintes: eles trabalhavam cada um em seu próprio computador, espalhados pela universidade e fora dela, conectados à infraestrutura sem partilhar uma bancada. O laboratório de IA que fui descrever materializa-se no computador individual de cada pesquisador e na infraestrutura computacional acessada à distância, e essa dispersão foi o primeiro achado do trabalho de campo, antes de ser um traço a contornar. A tese inteira se organizou em torno do trabalho de seguir esse laboratório para onde ele de fato estava, e não para onde eu supunha que estivesse.

Chamei este trabalho de uma tecno-etnografia, e a grafia carrega um argumento que percorre os cinco capítulos. No \texttt{LaTeX} em que escrevo a tese, e que aprendi a usar no próprio campo ao ver os pesquisadores do \texttt{C4AI} trabalharem com ele, as chaves delimitam um bloco que pertence junto como unidade, ao passo que os colchetes delimitam o que é opcional e acrescentado de fora. Um método aplicado de fora a um objeto que lhe preexistisse seria o colchete; o conceito da prática que proponho é a chave, em que técnica, tecnologia e etnografia formam um só escopo. Não tomei a tecno-etnografia como um método pronto que eu traria para aplicar ao \texttt{C4AI}, e sim como um conceito que se constituiu no próprio gesto de pôr em relação o campo, a teoria, as minhas memórias de pesquisa e os modelos de linguagem com que escrevi, no sentido de relação que desenvolvi no capítulo~\ref{capitulo2} a partir de \textcite{strathern1995relation}. Aprender a escrever em \texttt{LaTeX} foi, ele mesmo, uma marca que o campo deixou em mim, no registro do ser afetado, e a notação com que inscrevo a tese é um vestígio do que observei e do que o campo fez comigo enquanto eu o observava.

Esse gesto de pôr em relação organizou os três movimentos analíticos que o capítulo~\ref{capitulo1} tornou explícitos e que cada capítulo seguinte retomou em escala e materialidade próprias. O primeiro foi o corte. Toda rede se daria a ver como infinita se eu não a delimitasse, e cortar, no sentido de \textcite{strathern2011cortando}, é a operação que a torna descritível, em vez de uma limitação que a empobrece de fora. Os quatro capítulos fazem quatro cortes, reflexivo no capítulo~\ref{capitulo1}, teórico no capítulo~\ref{capitulo2}, pela extensão da rede no capítulo~\ref{capitulo3} e pela densidade dela no capítulo~\ref{capitulo4}, e cada corte performa o \texttt{C4AI} como um objeto etnográfico distinto, no sentido de \textit{enactment} que \textcite{mol2002body} desenvolve, sem que os quatro se subsumam num só. O segundo movimento foi descrever as existências parciais que se compõem na pesquisa. Parti do princípio de simetria de Latour e o desdobrei até as existências parciais de \textcite{strathern2004partial}, entidades conectadas a mundos distintos sem serem redutíveis a nenhum deles, e foi assim que pude nomear os sistemas de inteligência artificial generativa como quase-máquinas e reconhecer o \texttt{claude} como objeto múltiplo, no sentido de \textcite{Law2004After}, existente de modo distinto em treinamento, em uso e em auditoria. O terceiro movimento foi a compostagem, o pôr a fermentar materiais que o desenho inicial da pesquisa mantinha separados: a regra de estudar a ciência em construção com a dispersão da pesquisa em IA, a simetria metodológica com a especificidade da IA generativa, a minha formação em ciências sociais com o letramento técnico que o campo me exigiu, e a separação prévia entre ciência e política com a evidência de que a pesquisa no \texttt{C4AI} é, ao mesmo tempo, decisão epistêmica e captação de recursos. 

Figurei o próprio método por uma imagem têxtil, a do \textit{patchwork}, e mais precisamente a do \textit{crazy quilt}, em que a heterogeneidade dos retalhos, a visibilidade das costuras e a recusa de um padrão geral correspondem ao que a tese precisou fazer ao compor materiais heterogêneos. Para mim, no sentido de Strathern e de Haraway, cortar é condição da descrição, e o \textit{crazy quilt} mostra as costuras em vez de escondê-las: foi assim que procurei escrever, deixando à vista as junções entre os materiais em lugar de simular uma superfície contínua. Esse modo de proceder se sustenta numa etnografia heterodoxa, no sentido de \textcite{Kofes2015Maconaria}: o método se ensaia enquanto se experimenta, e não se aplica pronto a um campo que o aguardasse. Foi por me deixar afetar pelo encontro com o campo, no registro do ser afetado, que cada conceito que convoquei foi chamado pelo que o campo exigia, e que a própria escrita da tese se deixou medir pelos mesmos instrumentos com que me comprometi a medir o \texttt{C4AI}.

É parte desses três movimentos o modo como escrevi esta tese junto com modelos de linguagem, o \texttt{claude} e o \texttt{claude code}. Não me aproximei deles como ferramentas neutras que apenas transmitiriam o que eu já sabia, e sim como mediadores que transformavam aquilo que atravessava, em coerência com o conceito de mediação técnica que retomei no capítulo~\ref{capitulo2} e com o argumento sobre agência distribuída que sustento desde o capítulo~\ref{capitulo1}. Fazer pesquisa não só sobre, mas com a inteligência artificial generativa foi a inflexão que distingui da etnografia digital, e a auditabilidade pública das técnicas, nos repositórios do \texttt{GitHub}, é a contraparte material e ética desse compromisso. A análise lexicométrica das figurações, no capítulo~\ref{capitulo2}, e o mapeamento bibliométrico do campo brasileiro foram desenvolvidos com o \texttt{claude code} a partir das minhas especificações, e li o que voltou transformado sem que isso desloque para o modelo a agência da análise. Fazer-com, aqui, é manter a mediação presente como aquilo que a tecno-etnografia tem por compromisso tornar visível, sem apagá-la nem cedê-la.

Por que descrevi a tecnociência por esses movimentos, e não por outros, é uma decisão que o capítulo~\ref{capitulo2} fundamentou e que a análise lexicométrica deu a ver. A tecnociência vem sendo descrita há quatro décadas num vocabulário militar-industrial, o de alistar e recrutar aliados, vencer controvérsias e submeter o adversário a provas de força, e descrever o \texttt{C4AI} e o \texttt{Spira} nesse vocabulário herdado significaria naturalizar uma figuração específica, a da tecnociência como combate, e fechar de antemão a possibilidade de descrevê-la de outros modos. A análise que fiz de seis obras de Latour mostrou que a densidade do rótulo militar é baixa em \emph{Laboratory Life}, salta em \emph{Science in Action} e permanece alta em \emph{Pandora's Hope}, e que, quando Latour reinterpreta a própria teoria nos artigos metateóricos, esse vocabulário recua e cedem lugar os rótulos da topologia e do têxtil. A rede de Latour é, ela própria, um tecido, fibrosa e feita de fios e laços tecidos, e por isso a escolhi pôr em relação com a figura do berço de gato de \textcite{Haraway2023Ficar}, em que o conhecimento é o jogo de passar um fio de uma mão a outra, que só persiste enquanto há quem o receba e devolva, sem \textit{front} e sem adversário. O exército de aliados descreve como a tecnociência funciona; a figura de barbante descreve como ela funciona e, no mesmo gesto, propõe que poderia funcionar de outro modo. Foi com esse segundo repertório, ao lado da aliança sem ferir de Stengers, que escolhi descrever o que segui, e a postura que sustentei foi a de não afundar na denúncia que só vê extração nem na celebração que só vê solução.

Onde cheguei, se é que cheguei a um lugar, é a descrições situadas de dois pontos da mesma rede, e não a uma tese causal sobre como funcionam as parcerias público-privadas em IA nem a um modelo que se pudesse aplicar a outros casos. No capítulo~\ref{capitulo3}, ao seguir Fábio Cozman e Cláudio Pinhanez, o laboratório se estendeu até as corporações, os contratos, os eventos para investidores e a trajetória de mais de um século da \texttt{IBM}, da máquina de tabulação de Hollerith no Censo dos Estados Unidos de 1890 aos ecossistemas contemporâneos de IA. O que pude descrever, no caso do \texttt{C4AI}, foi um arranjo em que o valor foi retido pelo controle de um ecossistema de inovação, em que a multiplicação de centros criou as próprias condições do encerramento, e em que a paridade contratual de 2016 inscreveu desde o início o baixo custo da saída corporativa. Tomei as falas de Fábio e Cláudio como enunciações situadas de quem fala da posição de porta-voz do arranjo, e não como explicação de como a relação entre a \texttt{IBM} e o \texttt{C4AI} de fato funcionava nem como prova de um padrão que governasse esse tipo de parceria. No capítulo~\ref{capitulo4}, ao seguir Marcelo Finger e o \texttt{Spira}, o laboratório se materializou na cadeia de inscrições que vai do som da voz hospitalar ao espectrograma e à classificação pela rede neural, e o que pude descrever foi a tradução de uma escuta clínica incorporada em inscrições transportáveis, e a circunscrição que a pandemia negociou de dentro, quando os modelos que atingiam mais de noventa e cinco por cento de acurácia na primeira fase caíram para menos de trinta e seis por cento na segunda. Os dois capítulos descrevem o mesmo laboratório de pontos diferentes da rede, um a partir do arranjo institucional que torna a pesquisa possível, outro a partir das práticas situadas em que ela se faz, sem que isso faça de uma rede maior ou menor que a outra.

Desses dois pontos emerge um traço comum que a tese sustenta sobre o que é fazer pesquisa em IA neste contexto. O \enquote{laboratório de IA} não corresponde às expectativas moldadas pela física ou pela química experimental clássica, com seus espaços delimitados e suas rotinas de bancada. O que descrevi foi uma tecnociência distribuída, infraestruturalmente dependente, organizacionalmente híbrida e financeiramente acoplada a ciclos corporativos globais, e a sala de máquinas que João me mostrou em 2022 é a figura material dessa condição. A inteligência artificial não começa no algoritmo: começa na extração mineral, no cristal de silício, na eletricidade que corre pelos transistores e na infraestrutura computacional que poucos centros no mundo possuem em escala. Ensinou-se pedras a fazer cálculos, e o uso cotidiano da IA tende a esconder essa cadeia de mediações que vai do mineral à nuvem. Seguir o \texttt{C4AI} e o \texttt{Spira} foi também seguir essa cadeia até onde ela toca a Terra, e reconhecer que a dependência da infraestrutura computacional, no Brasil, é dependência de quem a fabrica e a comanda.

A coexistência de duas figurações atravessa a tese e foi o que me permitiu chegar até aqui sem reduzir um registro ao outro. Penso e construo o argumento por uma figuração têxtil, feita de fios, tramas, costuras, nós e cortes, que trabalha no tempo e na espessura do pensamento que se tece e em que o ato de relacionar nunca vem dado de antemão. Ao mesmo tempo, inscrevo o percurso da tese em diagramas e mapas que põem em relação numa superfície plana, na simultaneidade do que se apreende de um só golpe de vista. Os mapas de translação do capítulo~\ref{capitulo3} e da cadeia de inscrições do capítulo~\ref{capitulo4} são a trama que se mostra; o trabalho de seguir, cortar e costurar de que eles resultam é a urdidura que se pratica. A figuração têxtil resiste à superfície, porque o que ela figura é o próprio ato de relacionar em seu fazer-se, e um mapa mostra os nós e as conexões, isto é, as coisas já postas em relação, mas não mostra o relacionar que as pôs assim. Pratica-se a urdidura, mostra-se a trama. Habito essa tensão, que não chega a ser contradição, e a deixo aberta: a coexistência das duas figurações, em vez da escolha entre elas, é parte do que proponho.

Para onde se pode ir é a parte que prefiro deixar em aberto, porque a rede que descrevi estava, no momento em que encerrei a escrita, em plena reconfiguração. Em dezembro de 2025, o encerramento do laboratório da \texttt{IBM} no Brasil e da parceria \texttt{C4AI-IBM} transformou esta tese no registro não planejado de uma trajetória completa, da constituição em 2020 à dissolução cinco anos depois, e deixou em jogo a redistribuição dos recursos, das competências e dos vínculos que o arranjo havia organizado por cinco anos. O que permanece em aberto não se fechou por falta de rastreamento, e sim porque toca no que corre por canais que não acompanhei, os cálculos internos, as negociações bilaterais e as deliberações estratégicas que a opacidade do núcleo decisório manteve fora do meu alcance. Foi o corte que fiz ao seguir os porta-vozes, no sentido de Strathern, que me deu a circulação institucional e me subtraiu a decisória, e essa opacidade é parte do que circula nessas cadeias e condição de seu funcionamento, e não uma falha da descrição. 

Há uma segunda condição que esta tese situa, e que não diz respeito à infraestrutura, mas ao pensamento. O mapeamento bibliométrico que fiz no capítulo~\ref{capitulo2}, sobre o Catálogo de Teses e Dissertações da \texttt{Capes}, a coleção brasileira da \texttt{SciELO} e a base internacional \texttt{OpenAlex}, registrou em forma quantificada uma emergência tardia e tematicamente deslocada dos estudos brasileiros sobre inteligência artificial nas ciências humanas e sociais, marginalidade que se confirma da pós-graduação ao periódico e do recorte nacional ao internacional. Interpreto esse resultado menos como ausência de pensamento sobre a tecnologia e mais como uma defasagem do vocabulário com que esses estudos delimitam a inteligência artificial e as tecnologias correlatas. Situo esta tese nesse campo em formação, e a defasagem que o mapeamento tornou legível é parte do problema que carrego para o fechamento: se a pesquisa em IA no Brasil depende da infraestrutura de quem a fabrica, ela depende também de um vocabulário crítico que as ciências humanas brasileiras ainda estão construindo para descrevê-la, e esta tese é uma contribuição a essa construção tanto quanto uma descrição do \texttt{C4AI} e do \texttt{Spira}.

O problema que a tese tornou visível, e ao qual não entrego uma conclusão, diz respeito ao lugar da universidade pública quando a fronteira da produção de conhecimento em IA migra para a infraestrutura corporativa. A sala de máquinas que João me mostrou em 2022 era o \texttt{Arandu}, com oito \texttt{GPUs}, expressivo no contexto brasileiro e modesto em escala global; em fevereiro de 2026, a \texttt{USP} inaugurou o \texttt{JAIRU}, com noventa e seis \texttt{GPUs} de arquitetura mais recente, ao custo de cerca de quarenta milhões de reais, e com ele as mesmas perguntas sobre quem financiará a manutenção, como garantir acesso equitativo a pesquisadores fora de São Paulo e como formar conhecimento sobre sistemas dessa complexidade sem depender de quem os fabrica. Se os vínculos com quem comanda a infraestrutura computacional se tornaram condição de funcionamento para a pesquisa brasileira em IA, a pergunta que ofereço a outras pesquisas é sob quais arranjos esses vínculos vêm sendo buscados, com quais salvaguardas e em direção a quais finalidades. A essa pergunta não entrego uma resposta; ofereço um problema que a descrição tornou visível, e o deixo à disposição de outros rastreamentos, em particular de um retorno ao campo que a continuação do \texttt{Spira-BM}, em curso para além do recorte desta tese, ainda pode tornar possível.

Resta dizer que esta tese não se conclui no sentido em que se fecha uma caixa-preta. Aprendi com Fábio, Cláudio e Marcelo, no sentido em que \textcite{Ingold2019Antropologia} entende a antropologia, como compromisso de levar a sério quem se toma por interlocutor, e foi esse aprendizado que me deu fôlego para escrever. Descrever criticamente um arranjo do qual eu também era parte, sem converter essas pessoas em réus, foi o que me exigiu o esforço que tomei de Stengers, o de fazer alianças sem ferir, e a forma que encontrei foi dirigir a descrição aos arranjos que governaram o vínculo, e não às pessoas que os ocuparam dentro das margens estreitas que esses arranjos deixavam. Cheguei ao limite do que a minha posição parcial e situada sustenta, no sentido de conhecimento situado de \textcite{Haraway1991Situated}, e descrever esse limite como tal, em vez de ultrapassá-lo, é parte do que entendo por responsabilidade.

Volto, para encerrar, à sala de vidro. A transparência seletiva daquele espaço, que permitia ver o interior sem deixar entrar, foi a primeira figura do que esta tese veio a descrever: uma pesquisa que se podia acompanhar de fora, atrás do vidro, mas cujo núcleo decisório permanecia atrás de uma porta que não se abria para mim. Descrevi o que o vidro deixava ver, segui os atores até onde me levaram, e nomeei a porta fechada como porta fechada, em vez de fingir que a havia atravessado. É isso o que a minha posição alcançou, e é até aí que vai a descrição que assumo como minha.